\documentclass[11pt]{article}
\usepackage[T1]{fontenc}\usepackage[utf8]{inputenc}
\usepackage{lmodern,microtype}
\usepackage[a4paper,margin=26mm]{geometry}
\setlength{\emergencystretch}{3em}
\usepackage{amsmath,amssymb,amsthm,booktabs,array,hyperref}
\hypersetup{colorlinks=true,linkcolor=blue,urlcolor=blue,citecolor=blue,pdftitle={A reciprocal-deformation proof candidate for Ni Conjecture 3.1},pdfauthor={Reiwa}}
\newtheorem{theorem}{Theorem}
\newtheorem{lemma}{Lemma}
\newtheorem{proposition}{Proposition}
\newcommand{\qp}[3]{(#1;#2)_{#3}}
\newcommand{\F}[2]{ {}_4\phi_3\!\left[\begin{matrix}#1\\#2\end{matrix};R,R\right] }
\title{A reciprocal-deformation proof candidate for\\Ni's Conjecture 3.1 on a truncated double sum}
\author{Reiwa --- Tournament round 6}
\date{16 November 2026}
\begin{document}
\maketitle
\begin{abstract}
We present a complete proof candidate for the original all-$n$ assertion of
Conjecture 3.1 in He-Xia Ni's \emph{Supercongruences Arising from Truncated
Appell Series}, arXiv:2607.04583v2. The proposed argument combines an exact
terminating zero already implicit in Ni's proof with a reciprocal symmetry
obtained from Sears' classical balanced transformation. The symmetry forces
vanishing of a transverse derivative, which is then converted into a double
cyclotomic zero. The argument includes composite $n$ and checks all local
regularity conditions. It is submitted for independent checking, not as an
independently verified discovery or a claim of priority. Exact finite tests
support the calculation but are not used as a proof.
\end{abstract}

\section{Original problem and scope}
For $k\geq0$ define
\[
\qp{a}{R}{k}=\prod_{r=0}^{k-1}(1-aR^r),\qquad
\qp{a_1,\ldots,a_s}{R}{k}=\prod_{t=1}^s\qp{a_t}{R}{k}.
\]
Let $n\equiv3\pmod4$ be a positive integer and set $m=(n-1)/2$.
Thus $n\geq3$ and $m$ is odd. The exact target is
\begin{equation}\label{eq:original}
 S_n(q)=\sum_{i=0}^{m}\sum_{j=0}^{m}
 \frac{\qp{q}{q^2}{i+j}\qp{q}{q^2}{i}\qp{q,-q}{q^2}{j}
 q^{2i+2j}}
 {\qp{q^2}{q^2}{i}^{\,2}\qp{q^2}{q^2}{j}^{\,2}
 \qp{-q^2}{q^2}{j}}
 \equiv0\pmod{\Phi_n(q)^2}.
\end{equation}
This is Conjecture 3.1 on page 13 of Ni's version 2, dated 10 July
2026~\cite{Ni}. In particular the denominator has \emph{squares of}
$\qp{q^2}{q^2}{i}$ and $\qp{q^2}{q^2}{j}$, not shifted factorials with indices
$2i$ and $2j$. The ranges are the full square $0\leq i,j\leq m$.

Ni's Theorem 1.2 proves the same sum modulo $\Phi_n(q)$, with value zero
when $n\equiv3\pmod4$. Ni's Theorem 1.3 proves a congruence modulo
$\Phi_n(q)^2$ for a \emph{different} double sum. Neither result is stated as
a proof of \eqref{eq:original} in the source. We do not claim to settle
Ni's separate Conjecture 5.1.

A congruence for a rational function here means that a numerator after
clearing a denominator coprime to $\Phi_n$ is divisible by the stated power
of $\Phi_n$. All denominators needed below have that property.

\begin{theorem}[Full assertion under examination]\label{thm:main}
For every positive integer $n\equiv3\pmod4$, the rational function in
\eqref{eq:original} is divisible by $\Phi_n(q)^2$ in the localization of
$\mathbb Q[q]$ at polynomials coprime to $\Phi_n(q)$.
\end{theorem}
The following sections give the proposed proof. There is no primality,
size, convergence, or unproved auxiliary-conjecture assumption.

\section{Finite identities and their domains}
We use basic hypergeometric notation only for terminating finite sums:
\[
 {}_{r+1}\phi_r\!\left[\begin{matrix}a_0,\ldots,a_r\\b_1,\ldots,b_r
 \end{matrix};R,R\right]
 =\sum_{k=0}^{M}\frac{\qp{a_0,\ldots,a_r}{R}{k}}
 {\qp{R,b_1,\ldots,b_r}{R}{k}}R^k,
 \qquad a_0=R^{-M}.
\]
This convention deliberately stops at $M$; it does not invoke an infinite
series at a root of unity. The following are classical rational identities.

The terminating $q$-Chu--Vandermonde formula is
\begin{equation}\label{eq:chu}
 \sum_{k=0}^{M}\frac{\qp{R^{-M},a}{R}{k}}{\qp{R,c}{R}{k}}R^k
 =a^M\frac{\qp{c/a}{R}{M}}{\qp{c}{R}{M}}.
\end{equation}
Andrews' terminating $q$-Watson formula, in its odd case, is
\begin{equation}\label{eq:watson}
 {}_4\phi_3\!\left[\begin{matrix}R^{-M},a^2R^{M+1},c,-c\\
 aR,-aR,c^2\end{matrix};R,R\right]=0\qquad(M\text{ odd}).
\end{equation}
Both identities are reproduced in Ni's equations (1.17) and (1.20).

The Sears balanced transformation in the form needed here is
\begin{align}\label{eq:sears}
 \F{R^{-M},A,B,C}{D,E,F}
 &=\left(\frac{BC}{D}\right)^M
 \frac{\qp{DE/(BC),DF/(BC)}{R}{M}}{\qp{E,F}{R}{M}}\notag\\
 &\quad\times
 \F{R^{-M},A,D/B,D/C}{D,DE/(BC),DF/(BC)},
 \qquad DEF=ABC R^{1-M}.
\end{align}
For an accessible primary proof see Gao--Guo, Theorem 5.1,
equation (5.1)~\cite{GG}; the balancing condition is stated there explicitly.
Their prefactor is
$A^M\qp{E/A,F/A}{R}{M}/\qp{E,F}{R}{M}$, with the last two denominator
parameters $AR^{1-M}/E$ and $AR^{1-M}/F$. The elementary identity
\begin{equation}\label{eq:invert}
 \qp{u}{R}{M}=(-u)^M R^{M(M-1)/2}
 \qp{R^{1-M}/u}{R}{M}
\end{equation}
and the balance condition convert that version to \eqref{eq:sears}; the
last two denominator parameters may be interchanged. Chu--Li~\cite{CL}
provide another primary statement of the terminating Sears transformation.

All three identities are first taken at generic nonzero parameters, with
no vanishing denominator. Clearing the finitely many denominators gives
polynomial identities. They therefore remain valid under any specialization
where the displayed finite expressions are regular. No restriction
$|R|<1$ is needed for this use. Every root-of-unity specialization below
has $M<n$, and we check its remaining denominators explicitly.

\section{An exact vanishing deformation}
Define the following rational function for the fixed odd integer $m$:
\begin{equation}\label{eq:T}
 T(q,x)=\sum_{i=0}^m\sum_{j=0}^m
 \frac{\qp{xq}{q^2}{i+j}\qp{q}{q^2}{i}\qp{q,-q}{q^2}{j}
 q^{2i+2j}}
 {\qp{q^2}{q^2}{i}^{\,2}\qp{q^2}{q^2}{j}^{\,2}
 \qp{-xq^2}{q^2}{j}}.
\end{equation}
In particular $T(q,1)=S_n(q)$ exactly.

\begin{lemma}\label{lem:zero}
For $n=2m+1$ and $m$ odd,
\begin{equation}\label{eq:exactzero}
 T(q,q^{-n})=0
\end{equation}
as a rational-function identity in $q$.
\end{lemma}
\begin{proof}
Put $Q=q^2$. At $x=q^{-n}$ the factor
$\qp{xq}{Q}{i+j}=\qp{Q^{-m}}{Q}{i+j}$ vanishes whenever $i+j>m$.
For a fixed $j\leq m$, factoring its first $j$ factors and applying
\eqref{eq:chu} with $M=m-j$, $a=q$, $c=Q$ gives
\[
 \sum_{i=0}^{m-j}\frac{\qp{Q^{-(m-j)},q}{Q}{i}}{\qp{Q}{Q}{i}^{\,2}}Q^i
 =q^{m-j}\frac{\qp{q}{Q}{m-j}}{\qp{Q}{Q}{m-j}}.
\]
The finite-product reversal \eqref{eq:invert} gives
\begin{equation}\label{eq:reverse}
 q^{m-j}\frac{\qp{q}{Q}{m-j}}{\qp{Q}{Q}{m-j}}
 =q^m\frac{\qp{q}{Q}{m}}{\qp{Q}{Q}{m}}
 \frac{\qp{Q^{-m}}{Q}{j}}{\qp{q^{1-2m}}{Q}{j}}.
\end{equation}
Consequently
\begin{equation}\label{eq:watsonuse}
 T(q,q^{-n})=q^m\frac{\qp{q}{Q}{m}}{\qp{Q}{Q}{m}}
 {}_4\phi_3\!\left[\begin{matrix}Q^{-m},Q^{-m},q,-q\\
 Q,q^{1-2m},-q^{1-2m}\end{matrix};Q,Q\right].
\end{equation}
Use \eqref{eq:watson} with base $Q$, $M=m$, $a=q^{-n}$ and $c=q$.
Indeed $a^2Q^{m+1}=Q^{-m}$, $aQ=q^{1-2m}$ and $c^2=Q$.
Since $m$ is odd, the last expression is zero.
\end{proof}

This exact vanishing calculation is the zero-order ingredient already
present in the proof of Ni's Theorem 1.2. We do not claim it as new.
The prefactor in \eqref{eq:watsonuse} is the exact rational prefactor
$q^m\qp{q}{Q}{m}/\qp{Q}{Q}{m}$; replacing it prematurely by a merely
congruent expression is unnecessary.

\section{Reciprocity at every primitive root}
Fix any primitive $n$th root of unity $\zeta$. Set $R=\zeta^2$. Since $n$ is
odd, $R$ has order $n$, and
\begin{equation}\label{eq:relations}
 R^{-m}=\zeta,\qquad \zeta^2=R,\qquad R^{1-m}=R\zeta.
\end{equation}
Let
\begin{equation}\label{eq:W}
 D(x)=\qp{-xR}{R}{m},\qquad W(x)=D(x)T(\zeta,x).
\end{equation}
Every $x$-dependent denominator $\qp{-xR}{R}{j}$ in $T(\zeta,x)$ divides
$D(x)$. All constant denominators are nonzero, as checked below. Thus
$W(x)\in\mathbb Q(\zeta)[x]$, with degree at most $2m$: its $(i,j)$ term
after this clearing has degree at most $(i+j)+(m-j)=m+i$.

\begin{lemma}[Root-of-unity reciprocal symmetry]\label{lem:reciprocal}
With the above notation, for every odd $n=2m+1$,
\begin{equation}\label{eq:reciprocal}
 W(x)=x^{2m}W(1/x).
\end{equation}
This statement does not require $m$ to be odd.
\end{lemma}
\begin{proof}
Factor $\qp{x\zeta}{R}{i+j}=\qp{x\zeta}{R}{j}
\qp{x\zeta R^j}{R}{i}$. Since $\zeta=R^{-m}$, apply \eqref{eq:chu} with
$M=m$, $a=x\zeta R^j$, $c=R$ to the inner $i$ sum:
\begin{equation}\label{eq:Ij}
 \sum_{i=0}^m\frac{\qp{\zeta,x\zeta R^j}{R}{i}}{\qp{R}{R}{i}^{\,2}}R^i
 =\frac{I_j(x)}{\qp{R}{R}{m}},\qquad
 I_j(x)=(x\zeta R^j)^m\qp{R^{1-j}/(x\zeta)}{R}{m}.
\end{equation}
For generic $x$, finite-product cancellation gives
\begin{equation}\label{eq:Iratio}
 \frac{I_j(x)}{I_0(x)}=
 \frac{\qp{x\zeta}{R}{j}}{\qp{xR}{R}{j}},\qquad
 I_0(x)=(-1)^m R^{m(m+1)/2}\qp{xR}{R}{m}.
\end{equation}
For explicit verification of the ratio use
\[
 \frac{\qp{aR^{-j}}{R}{m}}{\qp{a}{R}{m}}
 =R^{-mj}\frac{\qp{R/a}{R}{j}}{\qp{R^{1-m}/a}{R}{j}},
 \qquad a=R/(x\zeta).
\]
Here $R/a=x\zeta$ and $R^{1-m}/a=x\zeta R^{-m}=xR$.
For $I_0$, the product equals $\prod_{r=1}^m(x\zeta-R^r)$; its roots can
be rewritten using $\zeta R^{-r}=R^{m+1-r}$, which proves the second
identity in \eqref{eq:Iratio}.

Substituting these formulas in the outer sum yields
\begin{align}\label{eq:Freduction}
 T(\zeta,x)&=\kappa\qp{xR}{R}{m}\mathcal F(x),\notag\\
 \kappa&=\frac{(-1)^m R^{m(m+1)/2}}{\qp{R}{R}{m}},\qquad
 \mathcal F(x)=\F{\zeta,-\zeta,x\zeta,x\zeta}{R,xR,-xR}.
\end{align}
The sum defining $\mathcal F$ terminates at $m$ because $\zeta=R^{-m}$.
Use \eqref{eq:sears} with
\[
 M=m,\quad A=-\zeta,\quad B=C=x\zeta,\quad
 D=R,\quad E=xR,\quad F=-xR.
\]
The balance condition is exact:
\[
 DEF=-x^2R^3,
 \qquad ABC R^{1-m}=-x^2\zeta^3R\zeta=-x^2R^3.
\]
The transformed parameters and prefactor are
\[
 BC/D=x^2,\quad D/B=D/C=\zeta/x,\quad
 DE/(BC)=R/x,\quad DF/(BC)=-R/x.
\]
It follows, as an identity in the indeterminate $x$, that
\begin{equation}\label{eq:Freciprocal}
 \mathcal F(x)=x^{2m}
 \frac{\qp{R/x,-R/x}{R}{m}}{\qp{xR,-xR}{R}{m}}\mathcal F(1/x).
\end{equation}
By \eqref{eq:W} and \eqref{eq:Freduction},
$W(x)=\kappa\qp{xR,-xR}{R}{m}\mathcal F(x)$.
Multiplying \eqref{eq:Freciprocal} by its prefactor proves
\eqref{eq:reciprocal}. Since $\deg W\leq2m$, both sides of
\eqref{eq:reciprocal} are polynomials, so the identity extends across any
exceptional $x$ excluded by the intermediate divisions. We do not assume
that the degree is exactly $2m$.
\end{proof}

\section{Regularity, derivatives, and the squared modulus}
We now check the specializations and complete the proposed proof of
Theorem~\ref{thm:main}.

For $1\leq r\leq m<n$, $1-R^r\ne0$ because $R$ has order $n$.
Also $1+R^r\ne0$: otherwise an element of the odd-order cyclic group
generated by $R$ would equal $-1$, which has order two. This works for
\emph{every} odd $n$, including composite $n$.
It proves that all denominators in $T(q,x)$ are nonzero at $(q,x)=(\zeta,1)$,
that $D(1)\ne0$, and that $\kappa$ is well defined. It also proves
regularity of both terminating Sears sums at $x=1$, including their
parameters $R/x$ and $-R/x$. During the use of Sears at generic $x$, any
additional zeros exclude only finitely many $x$; the resulting identity
is an identity of rational functions before polynomial clearing.

For completeness, the extra denominators in \eqref{eq:watsonuse} are
regular at $q=\zeta$ as well. The exponents in
$\qp{q^{1-2m}}{q^2}{m}$ are $1-2m+2r=2-n+2r$, $0\leq r<m$;
modulo $n$ they are $2,4,\ldots,2m$. Neither these minus-factors nor the
corresponding plus-factors vanish. All are nonzero rational functions
at generic $q$. Thus the exact zero in Lemma~\ref{lem:zero} is legitimately
specialized at $(\zeta,1)$.

By Lemma~\ref{lem:zero}, $T(\zeta,1)=0$, and therefore $W(1)=0$.
Differentiate the polynomial identity \eqref{eq:reciprocal}:
\[
 W'(x)=2m x^{2m-1}W(1/x)-x^{2m-2}W'(1/x).
\]
At $x=1$ this says $2W'(1)=2mW(1)=0$. In characteristic zero,
$W'(1)=0$. Differentiating \eqref{eq:W} and using $D(1)\ne0$ now gives
\begin{equation}\label{eq:Txzero}
 T_x(\zeta,1)=0.
\end{equation}
This step uses \emph{both} the reciprocal symmetry and the known zero at
$x=1$. Reciprocal symmetry alone does not assert a double zero.

The rational function $T(q,x)$ is jointly holomorphic in a neighborhood
of $(\zeta,1)$, by the denominator check just given. The curve
$x=q^{-n}$ is holomorphic near the nonzero point $q=\zeta$ and passes
through $(\zeta,1)$. We may therefore differentiate the exact rational
identity $T(q,q^{-n})=0$ there:
\begin{equation}\label{eq:chain}
 0=T_q(\zeta,1)-n\zeta^{-n-1}T_x(\zeta,1)
   =T_q(\zeta,1)-n\zeta^{-1}T_x(\zeta,1).
\end{equation}
Equation~\eqref{eq:Txzero} implies $T_q(\zeta,1)=0$.
Since $S_n(q)=T(q,1)$,
\[
 S_n(\zeta)=0,\qquad S_n'(\zeta)=0.
\]
There are no infinite sums, limiting integrals, or exchanges of infinite
operations: all differentiations are of finite rational expressions.

Finally put
\[
 C_m(q)=\qp{q^2}{q^2}{m},\qquad E_m(q)=\qp{-q^2}{q^2}{m},\qquad
 N_n(q)=C_m(q)^4E_m(q)S_n(q)\in\mathbb Z[q].
\]
Every summand's denominator divides this common denominator. The
nonvanishing checks prove $C_m(\zeta)E_m(\zeta)\ne0$ at every primitive
$n$th root. Thus $N_n(\zeta)=N_n'(\zeta)=0$ at each such root.
The cyclotomic polynomial has distinct roots in characteristic zero;
therefore $\Phi_n(q)^2$ divides $N_n(q)$ in $\mathbb C[q]$, and hence in
$\mathbb Q[q]$ by monic polynomial division. Since $N_n$ and $\Phi_n$ have
integer coefficients and $\Phi_n^2$ is monic, the quotient is in
$\mathbb Z[q]$. This proves precisely the asserted rational congruence.
\hfill$\square$

\section{Exact checks and limitations of computational evidence}
The proof candidate does not rely on a finite search. The accompanying
scripts provide independent checks of its algebraic predictions.
\begin{itemize}
\item \texttt{exact\_check.py} computes the common numerator $N_n$ in
$\mathbb Z[q]/(\Phi_n^3)$ directly from the original double sum, avoiding
floating-point roots and denominator inversion. For
$n=3,7,11,15,19,23,27,31,35,39,43,47$, its remainder modulo $\Phi_n^2$
is exactly zero. For $n\geq7$ in this list the remainder modulo
$\Phi_n^3$ is nonzero, so the order is exactly two. The cases $15,27,35,39$
are composite. These assertions and the residual polynomials are saved in
\texttt{exact\_checks.json}.
\item Direct symbolic simplification gives
\[
 S_3(q)=\frac{(q^2+q+1)^3}{(q+1)^3(q^2+1)}.
\]
Thus the smallest case has order three, not just two; the theorem does
not claim uniform exact order two.
\item \texttt{symmetry\_explore.py} clears the $x$ denominators in the
original double sum and reduces coefficients modulo $\Phi_n(q)$.
It checks the reciprocal coefficient symmetry exactly for $n=3,7,11$.
For $n=3$ its common-numerator polynomial is $3q(x-1)^2$ modulo $\Phi_3$.
The additional factor compared with $W$ is constant in $x$, so it does
not change reciprocal symmetry.
\end{itemize}
A successful finite test is not evidence that an unexamined specialization
is automatically legal; that is why the generic identities, balance,
termination, and regularity checks appear in the proof itself.

\section{Source, status, and novelty audit}
The source document and the two primary papers containing Sears
transformations were retrieved and read locally. The source's displayed
Conjecture 3.1 was checked against its rendered page 13, not reconstructed
from a secondary summary. The original is the all-positive-integer
$n\equiv3\pmod4$ assertion in \eqref{eq:original}.

The finite identities used here are classical, and Ni already uses the
Chu--Vandermonde and Watson zero. The additional proposed step is applying
Sears at a primitive root to obtain the reciprocal polynomial in the
\emph{same deformation parameter}, and then using its stationary zero
and the curve $x=q^{-n}$ to lift the modulus. This is a new application
within this submission, not a new general transformation. Whether that
meets a requirement for a substantively new method is a separate
judgment; the author does not certify that gate by assertion.

As of 16 November 2026, bounded searches for the paper's title, arXiv
identifier, and combinations of Ni, Conjecture 3.1, Appell, Sears, and
proof did not locate an existing proof of this exact conjecture. The
latest original version located in that search was version 2. This is
not a proof that no proof exists, and no claim of priority is made.
The corpus audit and retrieval limits are recorded in the accompanying
\texttt{audit.txt}. A candidate proof can be mathematically complete while
its novelty and open-status audit remain incomplete.

\begin{center}\small
\begin{tabular}{@{}p{0.09\linewidth}p{0.84\linewidth}@{}}\toprule
Gate & Assessment at submission\\\midrule
G1 & Exact existing named conjecture: Ni, Conjecture 3.1.\\
G2 & Original v2 read; full square ranges, factors, and squared modulus retained.\\
G3 & Conjectural in original; bounded later search found no resolution. Not exhaustively certified.\\
G4 & A full all-$n$ proof candidate is given above, with no conditional hypothesis; independent review pending.\\
G5 & Classical identities plus the reciprocal-deformation lift. Substantive novelty is not independently established.\\
G6 & Original proof compared; public search and supplied corpus examined with explicit limits. No priority claim.\\\bottomrule
\end{tabular}
\end{center}

\begin{thebibliography}{9}
\bibitem{Ni} H.-X. Ni, \emph{Supercongruences Arising from Truncated Appell
Series}, arXiv:2607.04583v2, 10 July 2026. Conjecture 3.1, p.~13;
Theorems 1.2 and 1.3; equations (1.17), (1.20).
\url{https://arxiv.org/abs/2607.04583v2}.
\bibitem{GG} F. Gao and V. J. W. Guo, \emph{Contiguous relations and
summation and transformation formulas for basic hypergeometric series},
arXiv:1304.5830v1, 22 April 2013. Theorem 5.1, equation (5.1), and its proof.
\url{https://arxiv.org/abs/1304.5830}.
\bibitem{CL} W. Chu and N. N. Li, \emph{Nonterminating extensions of the
Sears transformation}, Sarajevo Journal of Mathematics \textbf{12}(25),
no.~2 (2016), 205--215, doi:10.5644/SJM.12.2.07.
The terminating identity is equation (1), p.~206.
\url{https://www.sjm.ba/index.php/sjm/article/download/189/187/189}.
\end{thebibliography}
\end{document}
